\section{Introduction}
\label{sec:intro}

Predictive modelling of athletic performance is one of the oldest quantitative
pursuits in science, and the statistical literature on sports is enormous: it
ranges from rating systems for chess and ball games to forecasting individual
records in running and swimming. Yet most sports-analytics research is
conducted on bespoke, private or rapidly-changing datasets, which makes
results hard to compare. Each paper tends to choose its own data split, its own
preprocessing and its own metric, so two reported ``state-of-the-art'' numbers
are frequently not measuring the same thing. \wcabench{} addresses this
problem in a setting that has been almost entirely overlooked by the machine
learning community: competitive speedcubing, as recorded by the World Cube
Association (\wca).

\subsection{Motivation}
\label{sec:motivation}

Three problems motivate the benchmark.
\paragraph{Fragmented research.}
Existing work on \wca{} data is scattered across single tasks. Prior studies
have predicted competition placement with kernel density estimation, fitted
linear models to world-record trajectories, and estimated human limits with
Gaussian processes and extreme-value theory. Because these studies use
different splits, metrics and preprocessing, their conclusions cannot be
compared.
\paragraph{No standardized evaluation.}
The existing ``CubeBench''-style benchmarks target \emph{cube solving}: they
evaluate whether a language model or agent can plan a sequence of moves to
restore a scrambled cube. Those benchmarks are valuable, but they are
unrelated to the \wca's real competition records and cannot answer questions
such as ``which model predicts a competitor's time more accurately?'' or
``which method estimates the probability of a DNF more reliably?''. There is,
to our knowledge, no standardized benchmark for prediction and inference on
real speedcubing data.
\paragraph{Domain-specific structure is ignored.}
\wca{} data has unusual structure that general-purpose pipelines do not
handle: longitudinal competitor histories that span years; cross-event skill
transfer across 17 active events; round formats (\code{ao5}, \code{mo3},
best-of-3) that change how an average is computed; special status codes
($-1$ for \DNF{}, $-2$ for \DNS{}, $0$ for no result); and a compound encoding
for multi-blind results. A model that ignores the ``average of 5'' trimming
rule, for instance, will systematically miscalibrate both time and placement
predictions. We argue that these constraints deserve first-class treatment in
a benchmark rather than ad-hoc handling per paper.

\subsection{Contributions}
\label{sec:contributions}

We introduce \wcabench, a benchmark for prediction and inference on \wca{}
competition data. Our contributions are:

\begin{enumerate}
  \item \textbf{A reproducible data infrastructure.} A columnar
        (Parquet-based) pipeline that decodes every domain-specific value,
        including the \wca{} multi-blind compound encoding, round-format
        averaging with outlier trimming, and special status codes
        (Section~\ref{sec:dataset}). The pipeline is single-command
        reproducible and ships with a data card that documents provenance,
        scale, splits, biases and permitted/forbidden uses.

  \item \textbf{Five formalized tasks with domain-specific challenges.} We
        define \task{1} result prediction (regression), \task{2} placement
        prediction (ranking), \task{3} DNF prediction (imbalanced
        classification), \task{4} human-limit estimation (extreme-value
        extrapolation) and \task{5} skill-transfer analysis (causal
        inference). Each task is specified by an input/output/metric triple
        together with a baseline taxonomy and the domain challenges that make
        it non-trivial (Section~\ref{sec:tasks}).

  \item \textbf{A leakage-safe evaluation protocol and stratified
        reporting framework.} We specify a rolling-window protocol with
        frozen statistics, a four-dimensional stratification scheme (event,
        skill level, time slice, region) and a statistical-testing toolbox
        with effect sizes, designed so that small apparent gains are not
        over-claimed (Section~\ref{sec:protocol}).

  \item \textbf{21 reference baselines with real measured numbers.} We report
        actual baseline results on the official export
        (Section~\ref{sec:results}), covering six methodological
        families---statistical, tree, deep, graph, Bayesian and causal---with at
        least four baselines per task. For example, tree-based regression
        attains \maelog{} $=0.164$ on \task{1}; the \wca{} psychology sheet, a
        Plackett--Luce model and a kernel-density simulation all attain Kendall
        $\kendall=0.767$ on \task{2}; and tree-based classification attains
        \aucpr{} $=0.373$ on \task{3}. We deliberately do not claim
        state-of-the-art accuracy: the benchmark's scientific value lies in the
        standardized question, not in a new model.
\end{enumerate}

\subsection{Findings at a glance}
\label{sec:findings}

Because we release real numbers rather than hypothetical ones, the paper also
contains a first set of empirical observations, reported both for a single
seed and as a mean over three seeds (Section~\ref{sec:results-multiseed}).
Simple, well-specified models remain surprisingly competitive. On \task{2}, a
Plackett--Luce model, a kernel-density simulation and the official psychology
sheet produce \emph{identical} Kendall $\kendall=0.767$ to four decimals in the
seed-42 run ($0.8132\pm0.0342$ across seeds), whereas a graph neural network
reaches only $0.567$ ($0.7440\pm0.1258$). On \task{1}, a sequence model (LSTM,
\maelog{} $=1.044$, or $0.9436\pm0.0754$ across seeds) is much worse than
tree-based regression (\maelog{} $=0.164$, or $0.1055\pm0.0417$); and on
\task{3} the difference between the best tree model ($\aucpr{} 0.373$) and the
historical-rate prior ($\aucpr{} 0.359$) is \emph{not} statistically
significant ($p=0.749$, Cohen's $d=0.002$), while a Bayesian Beta--Binomial
model with no tree machinery is in fact best on average
($0.2938\pm0.0570$). Together these results suggest that, at this sample size,
rule-aware features combined with tree models dominate generic deep and graph
architectures. For human-limit estimation, the exponential-decay fit yields a
three-by-three limit of $3.74$ seconds, while the change-point model gives
$3.02$ seconds and the most stable method (Gaussian process plus extreme-value
theory) gives $3.29$ seconds---all well above the commonly cited domain anchor
of $2.37$ seconds, a concrete illustration of extrapolation risk. Finally, on
skill transfer, the number of ``identifiable'' event pairs falls monotonically
as identification assumptions tighten: Spearman correlation marks $413$, causal
forest $312$, instrumental variables $270$, and difference-in-differences only
$6$ ($2$ on average across seeds), quantifying how much the naive approach
over-claims. These observations are discussed in Section~\ref{sec:discussion}.

\subsection{Ethics and scope}
\label{sec:scope}

\wcabench{} is an evaluation benchmark, not a prediction product. It uses only
the publicly released \wca{} database export and never scrapes private data. We
explicitly forbid gambling, betting and any use that targets, ranks or
stigmatizes individual competitors, and we provide an opt-out mechanism
(Section~\ref{sec:limitations}). Out of scope are cube-solving/step-optimization
tasks (covered by cube-solving benchmarks such as \citep{cubebench, cube_bench_mllm}), LLM spatial-reasoning
evaluation, and live broadcast prediction systems.
